\emph{Setting and inputs.} The fields and the unit are those of Definition \claimref{obj.field-theory.single-algebra-field}; the loop conventions are those of Proposition \claimref{synthesis.one-loop-qed-in-the-spacetime-algebra}: the Dirac matrices are the vectors of $\Cl(1,3)$, the units of the Feynman rules and of the Wick rotation are coefficients collected in front of each Clifford product (clause (5) there: on a massive line the unit cannot act inside the product), and the trace is $\mathrm{Tr}\,X = 4\langle X\rangle_{0}$ (clause (1), checked there and here against Dirac matrices over $\mathbb{Q}(i)$). The chirality element of the matrix reading is $\gamma_{5} = iI_{4}$, $i$ the coefficient; then $\gamma_{5}^{2} = -I_{4}^{2} = 1$ and $\gamma_{5}$ anticommutes with every vector, and in the Dirac representation it is the usual off-diagonal identity (checked). Imported (Fact \claimref{synthesis.imported-one-loop-integrals-and-pauli-villars-regulators}): Feynman parametrisation, the Wick-rotated integral $\int\frac{d^{4}\ell}{(2\pi)^{4}}(\ell^{2} - \Delta)^{-3} = -i/(32\pi^{2}\Delta)$, and Pauli-Villars regulator fields with weights of either sign. Every Clifford identity below is exact and is checked over the Gaussian rationals at rational momenta and masses.

\emph{(i) Trace.} Every blade but the scalar is traceless, so $\mathrm{Tr}(\gamma_{5}X) = 4i\langle I_{4}X\rangle_{0}$. For four vectors only the grade-four part of $abcd$ contributes, $a\wedge b\wedge c\wedge d = \lambda I_{4}$, and $\langle I_{4}\lambda I_{4}\rangle_{0} = -\lambda$; so $\mathrm{Tr}(\gamma_{5}abcd) = -4i\lambda$, totally antisymmetric because the outer product is (checked on two orderings of four sample vectors and on the frame, $\langle I_{4}\gamma_{0}\gamma_{1}\gamma_{2}\gamma_{3}\rangle_{0} = -1$).

\emph{(ii) Vector Ward identities.} $S(p)$ is the inverse of $p - M$, so $S(a)hS(b) = S(b) - S(a)$ whenever $a - b = h$ (the resolvent identity of the cited result, clause (4)). In $D_{1}$ the vertex $\gamma^{\nu}$ contracted with $k_{1}$ sits between $S(\ell)$ and $S(\ell - k_{1})$; in $D_{2}$ it sits between $S(\ell - k_{2})$ and $S(\ell - q)$, whose momenta also differ by $k_{1}$. Hence $k_{1\nu}D_{1} = \mathrm{Tr}[\gamma^{\mu}\gamma_{5}S(\ell - k_{1})\gamma^{\rho}S(\ell - q)] - \mathrm{Tr}[\gamma^{\mu}\gamma_{5}S(\ell)\gamma^{\rho}S(\ell - q)]$ and $k_{1\nu}D_{2} = \mathrm{Tr}[\gamma^{\mu}\gamma_{5}S(\ell)\gamma^{\rho}S(\ell - q)] - \mathrm{Tr}[\gamma^{\mu}\gamma_{5}S(\ell)\gamma^{\rho}S(\ell - k_{2})]$; the middle terms cancel and the sum is $B(\ell - k_{1}) - B(\ell)$ (checked for all $\mu$, $\rho$ at three masses, including $0$, and three loop momenta; the $k_{2}$ contraction likewise). The two-point numerator $\mathrm{Tr}[\gamma^{\mu}\gamma_{5}(a + M)\gamma^{\rho}(b + M)]$ has pseudoscalar part from the term $a\gamma^{\rho}b$ only (the $M^{2}$ term is a bivector and the $M$ terms are odd), so it is mass independent and depends on $a$, $b$ through $a\wedge b$; with $a = c + u$, $b = c + u - k$ about the centre $c = k/2$ it is odd in $u$ (checked). For the regulated sum, put $\ell = tu$: the numerator of $\sum_{i}c_{i}/((t^{2}Q(u) - M_{i}^{2})(Q(tu - k) - M_{i}^{2}))$ over the common denominator has degree at most $\deg - 6$ exactly when $\sum_{i}c_{i} = 0$ (checked for two and three fields, with the control $\sum c_{i} \neq 0$, where the degree is $\deg - 4$). So $\sum_{i}c_{i}B_{i} = O(|\ell|^{-5})$, the regulated integral converges absolutely, it is invariant under the shift $\ell \to \ell - k_{1}$, and $k_{1\nu}T^{\mu\nu\rho} = 0$. The same decay ($\sum c_{i} = 0$ removes the $|\ell|^{-3}$ term of $D_{1} + D_{2}$) makes $T$ itself finite.

\emph{(iii) Axial divergence.} Since $\gamma_{5}$ anticommutes with vectors, $q\gamma_{5} = -S(\ell - q)^{-1}\gamma_{5} - \gamma_{5}S(\ell)^{-1} - 2M\gamma_{5}$. Inserting this between $S(\ell - q)$ and $S(\ell)$ (cyclically) gives $q_{\mu}(D_{1} + D_{2}) = -2M(P_{1} + P_{2})$ plus the four two-point terms $-\mathrm{Tr}[\gamma_{5}S(\ell)\gamma^{\nu}S(\ell - k_{1})\gamma^{\rho}]$, $-\mathrm{Tr}[\gamma_{5}\gamma^{\nu}S(\ell - k_{1})\gamma^{\rho}S(\ell - q)]$ and their images under $(\nu, k_{1}) \leftrightarrow (\rho, k_{2})$ (checked as an identity for all $\nu$, $\rho$ at two masses and loop momenta). Each two-point term is of the kind in (ii): under the regulators it converges and is odd about its centre, so it integrates to zero. In $P_{1}$ the numerator $\mathrm{Tr}[\gamma_{5}(\ell + M)\gamma^{\nu}(\ell - k_{1} + M)\gamma^{\rho}(\ell - q + M)]$ has a pseudoscalar part only at first order in $M$ (five vectors at order $0$ and three at order $2$ are odd; at order $3$, $\gamma^{\nu}\gamma^{\rho}$ has none); at first order the terms with two factors $\ell$ vanish ($\ell\wedge\ell = 0$) and the terms linear in $\ell$ cancel in pairs, leaving $M\cdot 4i\langle I_{4}\gamma^{\nu}k_{1}\gamma^{\rho}k_{2}\rangle_{0}$, independent of $\ell$; $P_{2}$ gives the same (checked for all $\nu$, $\rho$, two masses and two loop momenta). No surface term arises: the integrand's numerator is constant. With two Feynman parameters, $2\int_{\mathrm{simplex}} = 1$, and $\Delta_{i} = M_{i}^{2} - \pi(x, y)$ with $\pi$ a quadratic form in $k_{1}$, $k_{2}$, the regulator $i$ contributes
\[
c_{i}(-2M_{i})\cdot M_{i}\cdot 8i\langle I_{4}\gamma^{\nu}k_{1}\gamma^{\rho}k_{2}\rangle_{0}\cdot\frac{-i}{32\pi^{2}M_{i}^{2}}\bigl(1 + O(M_{i}^{-2})\bigr) = -\frac{c_{i}}{2\pi^{2}}\langle I_{4}\gamma^{\nu}k_{1}\gamma^{\rho}k_{2}\rangle_{0}\bigl(1 + O(M_{i}^{-2})\bigr),
\]
and the physical field contributes $-2m(P_{1} + P_{2})$, zero at $m = 0$. Summing, with $\sum_{i \geq 1}c_{i} = -1$, gives the stated limit; the regulator masses and the individual weights drop out (checked: three weight sets with $\sum c_{i} = 0$ give $\tfrac12$; the rational part $\tfrac{1}{32}$ of the integral is recomputed from the radial integral $\int_{0}^{\infty}r^{3}(r^{2} + \Delta)^{-3}dr = 1/(4\Delta)$, whose antiderivative is checked exactly, the angular volume $2\pi^{2}$ and $(2\pi)^{-4}$). The overall Feynman-rule factors of the amplitude ($-1$ for the loop, the units of propagators and vertices) multiply $T$ and are not part of the statement. The magnitude $1/(2\pi^{2})$ is the standard one of Adler and of Bell and Jackiw (Fact \claimref{synthesis.imported-adler-bell-jackiw-axial-anomaly}, stated there in the index reading of the sources as $(\alpha/4\pi)\epsilon FF$; Peskin and Schroeder, ch. 19, cited from memory). That it does not depend on the regulator is not proved here: the regulator-free extension result (\claimref{synthesis.one-loop-vacuum-polarisation-extends-without-a-regulator-and-ward-preserving-extensions-differ-only-by-the-renormalisation-constants}, clause (iii)) shows that every covariant extension conserving both vector currents has the same axial divergence.

\emph{(iv) Field form.} $\langle I_{4}a_{1}k_{1}a_{2}k_{2}\rangle_{0}$ sees only $a_{1}\wedge k_{1}\wedge a_{2}\wedge k_{2}$, which is $\langle (a_{1}\wedge k_{1})(a_{2}\wedge k_{2})\rangle_{4}$; so it equals $\langle I_{4}f_{1}f_{2}\rangle_{0}$, and swapping the photons leaves it unchanged since $f_{1}f_{2}$ and $f_{2}f_{1}$ have the same grade-four part (checked). For $F = E + I_{4}B$ with $E = E^{k}\gamma_{k}\gamma_{0}$, $B = B^{k}\gamma_{k}\gamma_{0}$, $F^{2} = E^{2} - B^{2} + 2(E\cdot B)I_{4}$, so $\langle I_{4}F^{2}\rangle_{0} = -2E\cdot B$ (checked at two field values). Reading the momentum-space divergence as a field equation gives the standard $|\partial_{\mu}j_{5}^{\mu}| = \tfrac{e^{2}}{2\pi^{2}}|E\cdot B|$ per Dirac fermion (argued, the usual Fourier translation; sign by convention).

\emph{(v) Chiral fields.} $P_{L}a = aP_{R}$ for a vector $a$ and $P_{R}\gamma^{\nu}P_{L} = \gamma^{\nu}P_{L}P_{L}$, so moving each projector to the right through the next vector flips it and the three collapse to one, $\mathrm{Tr}[\gamma^{\mu}P_{L}a\gamma^{\nu}P_{L}b\gamma^{\rho}P_{L}c] = \mathrm{Tr}[\gamma^{\mu}a\gamma^{\nu}b\gamma^{\rho}cP_{R}]$, which is half the vector trace plus half the $\gamma_{5}$ trace (checked on four index triples at a rational loop momentum). With group matrices at the vertices, the two orientations of the loop carry $\mathrm{tr}(T_{a}T_{b}T_{c})$ and $\mathrm{tr}(T_{a}T_{c}T_{b})$, and Bose symmetry of the external legs leaves the axial part proportional to $d_{abc} = \mathrm{tr}(T_{a}\{T_{b}, T_{c}\})$ (argued, not checked: the standard reduction, cited from memory). Clause (v) is therefore proved for the projector identity and argued for the reduction to $d$. Pauli-Villars regulators of a Weyl field are Dirac fields (Fact, clause (iii)), and the anomaly is their surviving mass term, as in (iii).

\emph{(vi) One generation.} The content is the particle half of Theorem \claimref{synthesis.anomaly-freedom-forces-slot-content}(ii) with $m = 1$, all left-handed. Represent $X = (A, B, y)$ with $A \in \mathfrak{sl}(3)$, $B \in \mathfrak{sl}(2)$, $y \in \C$ on each summand $(r_{3}, r_{2})_{Y}$ as $A \otimes 1 + 1 \otimes B + yY$, with $-A^{\mathrm{T}}$ on $\bar 3$: sixteen dimensions. The tensor $d$ is symmetric, so it vanishes exactly when its components on a basis vanish; on the basis of eight elementary and diagonal matrices of $\mathfrak{sl}(3)$, three of $\mathfrak{sl}(2)$ and one hypercharge, all 364 symmetric components are zero, and every basis element has trace zero (checked). For the hypercharge the last statement is $\mathrm{tr}\,Y = 0$, the mixed gravitational condition of the imported anomaly coefficients; it is a trace count on the content, not a computed gauge-gravitational triangle. This contains the four conditions of the stored theorem and the $[SU(3)]^{3}$ coefficient; the remaining components vanish because $\mathfrak{sl}(2)$ has no cubic invariant and the components with a single traceless factor vanish. A single field does not cancel ($\mathrm{tr}\,Y^{3} = 6/216$ on the quark doublet, checked), so the cancellation is between the fields of one generation and is independent of the number of generations. $\square$
